\PassOptionsToPackage{hidelinks}{hyperref}
\documentclass[
  aps,
  prd,
  reprint,
  superscriptaddress,
  nofootinbib,
  floatfix,
  longbibliography
]{revtex4-2}

\usepackage[T1]{fontenc}
\usepackage[utf8]{inputenc}
\usepackage{amsmath,amssymb,amsthm,mathtools,bm}
\usepackage{booktabs,array}
\usepackage{microtype}
\IfFileExists{orcidlink.sty}{\usepackage{orcidlink}}{\newcommand{\orcidlink}[1]{}}

\newtheorem{definition}{Definition}
\newtheorem{proposition}{Proposition}
\newtheorem{criterion}{Criterion}
\newtheorem{remark}{Remark}

\newcommand{\Sthree}{S^3}
\newcommand{\D}{\mathcal D}
\newcommand{\vis}{\mathrm{vis}}
\newcommand{\crit}{\mathrm{crit}}
\newcommand{\eff}{\mathrm{eff}}
\newcommand{\TT}{\mathrm{TT}}
\newcommand{\BY}{\mathrm{BY}}
\newcommand{\MS}{\mathrm{MS}}
\newcommand{\bigO}{\mathcal O}
\newcommand{\dd}{\mathrm d}
\newcommand{\nabl}{\nabla}
\newcommand{\Inv}{\mathcal I}

\begin{document}

\title{Benchmark residual invariants in compact \texorpdfstring{\(S^3\)}{S3} cosmology}

\author{Andrei~Preece\,\orcidlink{0009-0004-5978-795X}}
\email{andre.preece@web.de}
\affiliation{Independent Researcher, Berlin, Germany}

\author{Boris~Batenin\,\orcidlink{0009-0000-7027-8583}}
\email{batenin@web.de}
\affiliation{Independent Researcher, Berlin, Germany}

\date{\today}

\begin{abstract}
We formulate a benchmark-grade, first-principles framework for compact
\(S^3\) cosmology.  The useful object is not a verbal explanation of the
dark sector, but a reproducible map from invariant inputs to data
likelihoods.  Starting from
\((M,g_{\mu\nu},\Sigma,U^\mu,T^{\rm vis}_{\mu\nu})\), with
\(M\simeq I\times S^3\), we define the residual stress tensor,
the proper-time background reconstruction, the \(S^3\) spectral
covariance kernels, the Noether-II and quasi-local constraints, and the
data map needed for CMB, BAO, supernova, growth, lensing, and topology
tests.  We then add the standard that makes the construction useful to
other groups: a benchmark manifest, units and priors for each invariant
parameter, a transfer-operator ledger, validation tests, posterior
predictive residuals, evidence comparisons, and public reference
outputs.  This submission defines the benchmark-specification version of
the framework; numerical reference spectra, kernels, and likelihood
tables are deliberately separated into a subsequent validation release.
The result is not a microphysical dark-sector model and does not assert
that topology alone explains cosmic acceleration.  It is a strict
invariant-to-data benchmark against which compact-curvature and
residual-stress hypotheses can be tested, reproduced, and falsified.
\end{abstract}

\maketitle

\section{Introduction}
\label{sec:introduction}

The problem is not to attach a new name to the dark sector.  The problem
is to state the object which a successful dark-sector theory must match.
In a relativistic theory that object must be tensorial, conserved under
the appropriate hypotheses, and reconstructible from observables without
choosing a coordinate gauge; the conventions are those of standard general relativity and relativistic cosmology~\cite{Wald1984,HawkingEllis1973,Peebles1993,DodelsonSchmidt2020,Weinberg2008}.

We use the following first-principles data:
\begin{equation}
  (M,g_{\mu\nu},\Sigma,U^\mu,T^{\vis}_{\mu\nu}),
  \qquad
  M\simeq I\times S^3,
  \qquad
  U^\mu U_\mu=-c^2 .
  \label{eq:first_data}
\end{equation}
Here \(g_{\mu\nu}\) is the spacetime metric, \(\Sigma\) is a compact
Cauchy slice diffeomorphic to \(S^3\), \(U^\mu\) is the comoving
timelike congruence, and \(T^{\vis}_{\mu\nu}\) is the specified visible
stress tensor, including baryons, photons, neutrinos, and any other
sector one chooses to count as visible.

The central definition is
\begin{equation}
  T^\D_{\mu\nu}:=\frac{c^4}{8\pi G}G_{\mu\nu}
  -T^{\vis}_{\mu\nu}.
  \label{eq:residual_tensor}
\end{equation}
It is a residual target tensor, not a microscopic model.  It is
``complete'' only in this restricted sense: once \(g_{\mu\nu}\) and
\(T^{\vis}_{\mu\nu}\) are specified, Eq.~\eqref{eq:residual_tensor}
determines the dark-sector target.  It does not determine whether the
target is produced by particles, fields, modified gravity written in
Einstein form, interactions, or effective coarse graining.

The present version adds the part that makes the framework predictive.
It gives a network of invariants, specifies which of them lead to
observables, and states the likelihood structure needed for CMB, BAO,
supernova, lensing, growth, topology, and quasi-local tests.  We keep
three rules throughout.  First, no quantity is called invariant merely
because it is useful.  Scalars, spectra, topological numbers,
covariants, and gauge-redundant objects are separated.  Second, any
claim that depends on a bridge assumption is marked conditional.  Third,
numerical claims are not promoted to predictions until a transfer
function, likelihood, and data covariance have been specified.

\section{Geometry and proper time}
\label{sec:geometry_time}

Let
\begin{equation}
  \dd s^2=-N^2(t)c^2\dd t^2+a^2(t)R_0^2\dd\Omega_3^2 ,
  \label{eq:metric_lapse}
\end{equation}
where \(\dd\Omega_3^2\) is the unit three-sphere metric.  The physical
curvature radius is
\begin{equation}
  A(t):=a(t)R_0 .
  \label{eq:curvature_radius}
\end{equation}
The coordinate \(t\) has no invariant meaning unless \(N\) is fixed.
Therefore we use
\begin{equation}
  f' := U^\mu\nabla_\mu f=\frac1N\frac{\dd f}{\dd t},
  \qquad
  H:=\frac{a'}a=\frac{\dot a}{Na}.
  \label{eq:proper_time}
\end{equation}
A dot denotes \(\dd/\dd t\); a prime denotes differentiation with
respect to proper time \(\tau\).\footnote{The dot is used only for
coordinate-time differentiation before lapse fixing; all dynamical
background equations below use the prime/proper-time convention.}

For \(S^3_A\),
\begin{equation}
  V=2\pi^2A^3,\qquad
  {\cal R}^{(3)}=\frac{6}{A^2},\qquad
  \operatorname{diam}(S^3_A)=\pi A .
  \label{eq:s3_geometry}
\end{equation}
The curvature parameter is
\begin{equation}
  \Omega_K=-\frac{c^2}{A_0^2H_0^2},
  \qquad
  A_0=\frac{c}{H_0\sqrt{-\Omega_K}}
  \quad(\Omega_K<0).
  \label{eq:omegaK_radius}
\end{equation}
Thus the flat limit is not a topological limit inside \(S^3\); it is
the geometric limit \(A_0H_0/c\to\infty\).

\section{Residual stress-energy}
\label{sec:residual}

\begin{proposition}[Covariance and conservation]
The tensor \(T^\D_{\mu\nu}\) defined in Eq.~\eqref{eq:residual_tensor}
is a symmetric rank-two tensor.  If the visible stress tensor is
separately conserved,
\begin{equation}
  \nabla_\mu T^{\vis\,\mu\nu}=0,
  \label{eq:visible_conserved}
\end{equation}
then
\begin{equation}
  \nabla_\mu T^{\D\,\mu\nu}=0 .
  \label{eq:dark_conserved}
\end{equation}
\end{proposition}

\begin{proof}
The Einstein tensor and \(T^{\vis}_{\mu\nu}\) are tensors, so their
linear combination is a tensor.  Taking a divergence of
Eq.~\eqref{eq:residual_tensor} gives
\begin{equation}
  \nabla_\mu T^{\D\,\mu\nu}
  =\frac{c^4}{8\pi G}\nabla_\mu G^{\mu\nu}
  -\nabla_\mu T^{\vis\,\mu\nu}.
\end{equation}
The first term vanishes by the contracted Bianchi identity; the second
vanishes by assumption.
\end{proof}

If visible and dark sectors exchange energy-momentum, the correct
statement is instead
\begin{equation}
  \nabla_\mu T^{\vis\,\mu\nu}=J^\nu,\qquad
  \nabla_\mu T^{\D\,\mu\nu}=-J^\nu ,
  \label{eq:interaction_current}
\end{equation}
so that the sum is conserved.  A reconstructed nonzero divergence is
therefore not a falsification of \(S^3\) topology.  It is evidence for
an inconsistent reconstruction, a nonconserved visible sector, or a
missing interaction current.

For the FLRW background define
\begin{equation}
  \epsilon_{\crit}:=\frac{3c^2H_0^2}{8\pi G},
  \qquad
  E(a):=\frac{H(a)}{H_0}.
\end{equation}
The Friedmann equations give
\begin{equation}
  \frac{\epsilon_{\rm tot}}{\epsilon_{\crit}}
  =
  E^2-\Omega_Ka^{-2},
  \label{eq:eps_total}
\end{equation}
and
\begin{equation}
  \frac{p_{\rm tot}}{\epsilon_{\crit}}
  =
  -\frac13
  \left[
  \frac{\dd E^2}{\dd\ln a}
  +3E^2-\Omega_Ka^{-2}
  \right].
  \label{eq:p_total}
\end{equation}
The proper-time acceleration equation and deceleration parameter are
\begin{equation}
  \frac{a''}{a}
  =
  -\frac{4\pi G}{3c^2}(\epsilon_{\rm tot}+3p_{\rm tot}),
  \qquad
  q=-\frac{aa''}{(a')^2}.
  \label{eq:acceleration_q}
\end{equation}

Let
\begin{equation}
  \Omega_{\vis}(a)=
  \Omega_ba^{-3}+\Omega_\gamma a^{-4}+\Omega_\nu(a).
\end{equation}
Then
\begin{equation}
  \frac{\epsilon_\D}{\epsilon_{\crit}}
  =
  E^2-\Omega_Ka^{-2}-\Omega_{\vis}(a),
  \label{eq:epsD}
\end{equation}
and
\begin{equation}
  \frac{p_\D}{\epsilon_{\crit}}
  =
  -\frac13
  \left[
  \frac{\dd E^2}{\dd\ln a}
  +3E^2-\Omega_Ka^{-2}
  \right]
  -\frac13\Omega_\gamma a^{-4}
  -P_\nu(a),
  \label{eq:pD}
\end{equation}
where \(P_\nu(a)=p_\nu(a)/\epsilon_{\crit}\).  These formulae are
background targets.  They do not split \(T^\D_{\mu\nu}\) uniquely into
clustering dark matter and smooth dark energy; that split requires
perturbative observables.

\section{\texorpdfstring{\(S^3\)}{S3} spectra and covariance kernels}
\label{sec:spectra}

The scalar Laplacian on \(S^3_A\) has eigenvalues
\begin{equation}
  -\Delta Y_{n\ell m}
  =
  \frac{n(n+2)}{A^2}Y_{n\ell m},
  \qquad
  n=0,1,2,\ldots ,
  \label{eq:scalar_spectrum}
\end{equation}
with degeneracy \((n+1)^2\) in the standard scalar-harmonic
normalization on three-spheres~\cite{RubinOrdonez1984,AbbottSchaefer1986,Jantzen1978,Higuchi1987,LindblomTaylorZhang2017}.  The first nonzero scalar gap is
\begin{equation}
  \lambda_1=\frac{3}{A^2}.
  \label{eq:spectral_gap}
\end{equation}
The addition kernel is
\begin{equation}
\begin{split}
  \Xi_{n,A}(x,y)
  &=
  \sum_{\ell,m}Y_{n\ell m}(x)Y^*_{n\ell m}(y),
  \\
  \Xi_{n,A}(0)&=\frac{(n+1)^2}{2\pi^2A^3}.
\end{split}
  \label{eq:addition_kernel}
\end{equation}
For a statistically isotropic scalar Gaussian field,
\begin{equation}
  C(x,y)=
  \sum_{n=0}^{\infty}A_n\,\Xi_{n,A}(x,y)
  \label{eq:covariance}
\end{equation}
is a covariance only if
\begin{equation}
  A_n\ge0,
  \qquad
  \sum_{n=0}^{\infty}A_n\,\Xi_{n,A}(0)<\infty .
  \label{eq:covariance_conditions}
\end{equation}
Stronger decay controls differentiability.  Equation
\eqref{eq:covariance_conditions} is a sharp admissibility condition:
without it the coefficients are not a covariance spectrum.

CMB and large-scale-structure observables enter through transfer
functions, as in the standard Boltzmann-code approach~\cite{CLASS2011,CLASSNonFlat2014,CAMB2000}.  Schematically, for CMB fields \(X,Y\in\{T,E,B,\phi\}\),
\begin{equation}
  C_\ell^{XY}(\theta)
  =
  \sum_{n\in{\cal N}_\ell}
  A_n(\theta)\,
  \Delta_{\ell n}^{X}(\theta)\,
  \Delta_{\ell n}^{Y}(\theta),
  \label{eq:closed_cl}
\end{equation}
where \({\cal N}_\ell\) is the allowed discrete mode set and
\(\Delta_{\ell n}^{X}\) are the closed-FLRW transfer functions with the
chosen normalization.  Equation \eqref{eq:closed_cl} is the predictive
replacement for any verbal statement about low-\(\ell\) anomalies.
Closed-FLRW hyperspherical transfer functions are already implemented in
non-flat Boltzmann technology; the compact benchmark adds the discrete
sum over allowed \(n\) and the covariance-admissibility condition in
Eq.~\eqref{eq:covariance_conditions}.  A deficit, oscillation,
alignment, or parity statistic is a derived likelihood statement, not a
prior conclusion.

Vector and tensor perturbations require their own operator ledger.
Scalar eigenvalues cannot be re-used for transverse vectors or
transverse-traceless tensors without the correct curvature shifts,
constraints, and gauge conditions.

\section{Noether-II, quasi-local, and gauge constraints}
\label{sec:noether}

Diffeomorphism invariance is governed by Noether's second theorem and the covariant phase-space identities, not
by ordinary global Noether-I charges~\cite{IyerWald1994,Wald1984}.
The quasi-local examples below use the Brown--York and Misner--Sharp constructions~\cite{BrownYork1993,MisnerSharp1964}; the quantum-clock caution uses the Page--Wootters benchmark~\cite{PageWootters1983}.  On a compact Cauchy slice without
boundary,
\begin{equation}
  \partial\Sigma=\varnothing ,
\end{equation}
the diffeomorphism Hamiltonian is a constraint for pure gauge
deformations.  Physical energy diagnostics arise from subregions
\(D\subset\Sigma\) with \(\partial D\ne\varnothing\), from matter
clocks, or from asymptotic/boundary structures when such structures
exist.

For a geodesic two-sphere of areal radius \(R_A\) in a time-symmetric
round \(S^3\) slice, the Brown--York energy with flat reference
subtraction is~\cite{BrownYork1993,York1972}
\begin{equation}
  E_{\BY}(R_A)
  =
  \frac{c^4R_A}{G}
  \left[
  1-\sqrt{1-\frac{R_A^2}{A^2}}
  \right],
  \qquad 0\le R_A\le A .
  \label{eq:by_energy}
\end{equation}
For \(R_A\ll A\),
\begin{equation}
  E_{\BY}(R_A)
  =
  \frac{c^4R_A^3}{2GA^2}
  +\bigO(R_A^5/A^4).
  \label{eq:by_small}
\end{equation}
The leading term has volume scaling.  A Newtonian interpretation
requires the constraint equations and a specified reference subtraction;
it is not automatic.

The Misner--Sharp energy in FLRW is
\begin{equation}
  E_{\MS}(R_A)=\frac{4\pi}{3}R_A^3\epsilon_{\rm tot}.
  \label{eq:ms_energy}
\end{equation}
Equations \eqref{eq:by_energy} and \eqref{eq:ms_energy} are different
quasi-local diagnostics: \(E_{\BY}\) uses an embedding/reference
subtraction, while \(E_{\MS}\) is the local spherical FLRW mass
function.  They should not be identified unless the geometric setting
and subtraction convention justify it.

For a gauge field on a compact slice, the Poisson equation has a
zero-mode constraint:
\begin{equation}
  \int_\Sigma \rho\,\dd V=0 .
  \label{eq:gauss_zero}
\end{equation}
Equivalently, the Green function must be zero-mode subtracted:
\begin{equation}
  -\Delta G_0(x,y)=\delta_y(x)-\frac1{V(\Sigma)} .
  \label{eq:zero_subtracted_green}
\end{equation}
This is a global solvability condition, not a claim that every local
flux integral vanishes.

\section{Predictive invariant-to-data pipeline}
\label{sec:pipeline}

A model becomes experimentally useful only after it supplies a map
\begin{equation}
  \theta\longmapsto m(\theta)
  \label{eq:model_map}
\end{equation}
from invariant parameters to a data vector.  A minimal data vector is
\begin{equation}
\begin{split}
  d=&\{C_\ell^{TT},C_\ell^{TE},C_\ell^{EE},C_\ell^{\phi\phi},
  D_M(z_i)/r_d,D_H(z_i)/r_d,\\
  &\mu_{\rm SN}(z_j),H(z_k),f\sigma_8(z_l),
  \xi_\pm(\vartheta_m),{\cal T}_{\rm top}\}.
\end{split}
\label{eq:data_vector}
\end{equation}
Here \({\cal T}_{\rm top}\) denotes any topology statistic, such as a
circle-search statistic for a multiply connected quotient
\(S^3/\Gamma\) or a geodesic-wrap statistic for simply connected
\(S^3\), in the sense of compact-topology searches, if the horizon
scale is large enough.

A corresponding parameter vector may be written
\begin{equation}
\begin{split}
  \theta=&\{A_0,H_0,\Omega_bh^2,\Omega_\gamma h^2,
  \Omega_\nu h^2,\\
  &A_n\ \text{or}\ P_{\cal R}(n),
  \alpha_{\D},\eta_{\rm int},\nu_{\rm nuis}\}.
\end{split}
\label{eq:parameter_vector}
\end{equation}
The symbols \(\alpha_\D\) parameterize the residual background or
perturbations, \(\eta_{\rm int}\) parameterizes possible interaction
currents \(J^\mu\), and \(\nu_{\rm nuis}\) denotes nuisance parameters
for foregrounds, calibration, galaxy bias, intrinsic alignments, and
selection functions.

The likelihood is
\begin{equation}
  -2\ln{\cal L}(\theta)
  =
  [d-m(\theta)]^T{\cal C}^{-1}
  [d-m(\theta)]
  +\ln\det{\cal C}
  +\mathrm{const.},
  \label{eq:likelihood}
\end{equation}
with a non-Gaussian replacement where the experiment requires it.  This
is the correct form of falsifiability.  One does not declare a fixed
curvature radius to be excluded in the abstract.  One reports a posterior
or Bayes factor after specifying data, priors, transfer functions, and
covariance.

\begin{criterion}[Transfer-grade predictive closure]
A compact \(S^3\) residual model is transfer-grade only if it supplies:
\begin{enumerate}
\item the metric and proper-time convention;
\item the visible stress tensor and any interaction current;
\item the scalar, vector, and tensor operator ledger;
\item admissible covariance spectra satisfying
      Eq.~\eqref{eq:covariance_conditions};
\item transfer functions for the data vector \eqref{eq:data_vector};
\item a likelihood and reproducible numerical implementation;
\item a falsification statement in posterior, Bayes-factor, or
      predictive-residual form.
\end{enumerate}
\end{criterion}

The practical implementation is a closed-FLRW Boltzmann pipeline with
discrete compact mode sums.  Existing non-flat Boltzmann solvers provide
the hyperspherical transfer-function technology.  Compact topology adds
the finite-volume mode ledger, possible quotient selection rules,
admissible covariance coefficients, and topology likelihoods, with compact-topology searches providing one reference implementation of the latter step~\cite{SahaCompact2025}.  A code should
therefore not be advertised as predicting an anomaly until it returns
the objects in Eq.~\eqref{eq:data_vector} with the covariance in
Eq.~\eqref{eq:likelihood}.


\section{Benchmark standard for scientific utility}
\label{sec:benchmark}

A framework is useful to physics only when another group can reproduce
its numerical claims without guessing conventions.  The benchmark object
is therefore
\begin{equation}
  {\mathfrak B}
  =
  ({\cal I},\Theta,\pi,{\cal T},{\cal N},{\cal V},{\cal R}),
  \label{eq:benchmark_object}
\end{equation}
where \({\cal I}\) is the invariant ledger, \(\Theta\) is the parameter
space with units, \(\pi\) is the prior measure, \({\cal T}\) is the
transfer map from invariants to observables, \({\cal N}\) is the noise
and covariance model, \({\cal V}\) is the validation suite, and
\({\cal R}\) is the public release manifest.  Omitting any one of these
entries turns a prediction into an interpretation.

\begin{definition}[Registered prediction]
A registered prediction is a tuple
\begin{equation}
  P=(\theta,\pi,{\cal T},d_{\rm blind},{\cal C},{\cal L},{\cal F})
  \label{eq:registered_prediction}
\end{equation}
fixed before the target data \(d_{\rm blind}\) are used.  The final
entry \({\cal F}\) is the failure criterion: posterior exclusion,
Bayes-factor loss, or posterior-predictive inconsistency.  A fit made
after seeing the data is not a prediction; it is a consistency test.
\end{definition}

\begin{proposition}[Coordinate-free reproducibility]
If the entries in Eq.~\eqref{eq:benchmark_object} are expressed through
scalar invariants, tensorial projections, spectra, or quasi-local
surface integrals, then the benchmark predictions are independent of
coordinate charts on \(M\) and on \(\Sigma\).
\end{proposition}

\begin{proof}
The data map \(m(\theta)\) uses only the invariant ledger and the
experiment's projection operators.  A coordinate change alters component
representations of \(g_{\mu\nu}\), \(T^{\cal D}_{\mu\nu}\), harmonic
bases, and surface embeddings, but not eigenvalues, multiplicities,
addition kernels, contractions with \(U^\mu\), likelihood scalars, or
surface integrals with the induced measure.  Hence the numerical output
is a function of geometric objects, not of labels.
\end{proof}


The minimal benchmark products are the following.
\begin{enumerate}
\item \emph{Invariant manifest}: definitions, units, domains, priors, and
exact/conditional status of all entries \(\Inv_i\).  This prevents hidden
convention changes and non-comparable analyses.
\item \emph{Transfer ledger}: the map from
\(\{\lambda_n,d_n,\Xi_{n,A},T^{\cal D}_{\mu\nu}\}\) to \(C_\ell\), BAO
distances, growth, lensing, and topology statistics.  This prevents
claims of prediction without a computable observable.
\item \emph{Validation suite}: analytic limits, including the flat
limit, large-radius limit, covariance positivity, \(\Lambda\)CDM
residual check, and small-sphere Brown--York expansion.  This prevents
numerical agreement that is only accidental.
\item \emph{Likelihood package}: data vector, covariance, masks,
nuisance parameters, priors, and posterior-predictive residuals.  This
prevents unfalsifiable comparison with selected plots.
\item \emph{Reference outputs}: versioned tables for spectra, kernels,
background residuals, and benchmark likelihood values.  This prevents
irreproducible code-dependent numbers.
\item \emph{Release manifest}: version hashes, exact datasets,
compilation/runtime environment, and license.  This prevents results
that cannot be audited or reused.
\end{enumerate}


The first numerical target is not a full explanation of CMB anomalies or
cosmological tensions.  The first target is more severe and more useful:
given a closed-FLRW background and a compact spectral ledger, reproduce
standard non-flat transfer-function limits, replace continuous long-mode
integrals by the correct \(S^3\) sums, and report the likelihood change
relative to the flat and simply-connected baselines.  Only after this
test is passed may one discuss low-\(\ell\) anomalies, matched circles,
or residual-stress interpretations.

Quasi-local observables obey the same rule.  A Brown--York or
Misner--Sharp map inferred from lensing is not evidence for topology by
itself.  It becomes useful only if the surface, subtraction convention,
mass definition, data projection, and covariance are fixed in the
manifest.  The output is then a geometric diagnostic that can agree with,
or fail against, the residual target.

\section{Observational uses}
\label{sec:observational}

The invariant framework produces several test classes.

First, background data constrain
\begin{equation}
  \epsilon_\D(a),\qquad p_\D(a),\qquad
  w_\D(a)=p_\D(a)/\epsilon_\D(a)
\end{equation}
through Eqs.~\eqref{eq:epsD}--\eqref{eq:pD}.  BAO, supernovae,
cosmic-chronometer, BBN, and CMB priors enter here through \(E(a)\)~\cite{Planck2018VI,Planck2018VII,PlanckPR4CamSpec,DESIDR2BAO2025,DESIDR2Lya2025,DESIDR2DE2025,PantheonPlus2022,Moresco2022,ACTDR62024},
\(r_d\), and the visible-sector densities.

Second, perturbations test the split between a clustering and a smooth
residual.  Distances alone cannot identify dark matter.  Lensing,
redshift-space distortions, CMB peaks, cluster abundance, and velocity
fields are needed.

Third, the discrete \(S^3\) spectra generate covariance predictions for
low multipoles, parity statistics, alignments, and cross-correlations.
These are not automatic explanations of anomalies.  They are likelihood
tests of Eq.~\eqref{eq:closed_cl} with admissible \(A_n\).

Fourth, compact topology can be tested by geodesic wrap-around
statistics on \(S^3\) and by matched-circle statistics on quotients
\(S^3/\Gamma\).  The absence of a circle signal constrains the quotient,
radius, orientation, and noise model used in that search; it is not by
itself a universal exclusion of closed positive curvature.

Fifth, quasi-local diagnostics such as \(E_{\BY}\) and \(E_{\MS}\) can
be compared with lensing mass maps and cluster observables only after
specifying the slicing, reference subtraction, mass definition, and
selection function.  This is a difficult but well-posed extension.

\section{Conditional extensions}
\label{sec:conditional}

Some useful mathematics is conditional rather than first-principles
physics.  If a hierarchy is defined by
\begin{equation}
  F_{n+2}=F_{n+1}+F_n ,
\end{equation}
then
\begin{equation}
  \frac{F_{n+1}}{F_n}\to
  \varphi=\frac{1+\sqrt5}{2}.
\end{equation}
This is a Perron--Frobenius fact about a recursion.  It is not a
spacetime prediction unless a separate bridge connects the recursion to
spectral or geometric data.

Similarly, graph spectra, Zeckendorf decompositions,
Born--Oppenheimer parameters, and quantum-cosmological phases are useful
only after their Hilbert spaces, operators, gauge fixing, and
observables are stated.  They may be compared with the residual target
\eqref{eq:residual_tensor}; they do not replace it.

\section{Discussion}
\label{sec:discussion}

The rigid content of the construction is
\begin{equation}
\begin{gathered}
  \text{compact }S^3
  +T^\D_{\mu\nu}
  +\text{ proper-time FLRW geometry}
  \\
  +\text{ discrete spectra}
  +\text{ Noether-II constraints}
  \\
  +\text{ compact gauge zero modes}
  +\text{ quasi-local diagnostics}.
\end{gathered}
\label{eq:core}
\end{equation}
This is enough to define a useful research programme.  It is not enough
to claim a dark-sector mechanism, an explanation of the Hubble tension,
or a proof of any famous conjecture.  The scientific value lies in the
opposite discipline: every extension must pass through invariant
definitions, transfer functions, and data likelihoods.

The natural journal standard for the cosmology part is Physical Review D:
the content is gravitation, cosmology, field theory, and data-facing
cosmological inference.  A Physical Review A version would have to be a
different paper, centered on quantum information, quantum measurement, or
atomic/optical analogues.  The present manuscript should therefore be
judged primarily as a PRD-style theory-and-phenomenology paper.


The standard for the next version is therefore not a larger list of
claims but a smaller set of benchmark numbers: reference spectra,
reference residual histories, reference covariance kernels, and
reference likelihood values.  These are the objects that let another
group reproduce, refute, or improve the construction.

\appendix

\section{Predictive invariant ledger}
\label{app:ledger}

This appendix records eighty-four entries.  A scalar invariant is a
number or scalar field.  A spectral invariant is an eigenvalue,
multiplicity, trace, or zeta quantity.  A natural covariant is tensorial
or geometric but not scalar.  A conditional entry requires an additional
bridge assumption.

\begin{enumerate}
\item \(\Inv_1=M\simeq I\times S^3\): spacetime product topology.
\item \(\Inv_2=\partial\Sigma=\varnothing\): absence of spatial boundary.
\item \(\Inv_3=\pi_1(S^3)=0\): simple connectivity.
\item \(\Inv_4=H_0(S^3)=\mathbb Z\): connectedness.
\item \(\Inv_5=H_1(S^3)=0\): no one-cycle homology.
\item \(\Inv_6=H_2(S^3)=0\): no two-cycle homology.
\item \(\Inv_7=H_3(S^3)=\mathbb Z\): orientation volume class.
\item \(\Inv_8=\chi(S^3)=0\): Euler characteristic.
\item \(\Inv_9=A(t)=a(t)R_0\): physical curvature radius.
\item \(\Inv_{10}=V=2\pi^2A^3\): spatial volume.
\item \(\Inv_{11}=\operatorname{diam}S^3_A=\pi A\): diameter.
\item \(\Inv_{12}={\cal R}^{(3)}=6/A^2\): scalar curvature.
\item \(\Inv_{13}=R^{(3)}_{ij}=2h_{ij}/A^2\): Ricci tensor.
\item \(\Inv_{14}=K=+1/A^2\): sectional curvature.
\item \(\Inv_{15}=\Omega_K=-c^2/(A_0^2H_0^2)\): dimensionless curvature.
\item \(\Inv_{16}=A_0H_0/c=1/\sqrt{-\Omega_K}\): curvature radius in Hubble units.
\item \(\Inv_{17}=U^\mu U_\mu=-c^2\): normalized congruence.
\item \(\Inv_{18}=f'=U^\mu\nabla_\mu f\): proper-time derivative.
\item \(\Inv_{19}=H=a'/a\): expansion rate.
\item \(\Inv_{20}=\Theta=\nabla_\mu U^\mu=3H\): expansion scalar.
\item \(\Inv_{21}=q=-aa''/(a')^2\): deceleration scalar.
\item \(\Inv_{22}=1+z=a_0/a\): redshift in FLRW.
\item \(\Inv_{23}=E(z)=H(z)/H_0\): dimensionless expansion.
\item \(\Inv_{24}=\chi(z)=\int_0^z c\,\dd z'/H(z')\): radial comoving distance.
\item \(\Inv_{25}=D_M(z)=A_0\sin[\chi(z)/A_0]\): transverse distance for \(S^3\).
\item \(\Inv_{26}=D_H(z)=c/H(z)\): Hubble distance.
\item \(\Inv_{27}=r_d\): drag-epoch sound horizon.
\item \(\Inv_{28}=D_M/r_d\): BAO transverse observable.
\item \(\Inv_{29}=D_H/r_d\): BAO radial observable.
\item \(\Inv_{30}=G_{\mu\nu}\): Einstein tensor natural covariant.
\item \(\Inv_{31}=\nabla_\mu G^{\mu\nu}=0\): contracted Bianchi identity.
\item \(\Inv_{32}=T^{\vis}_{\mu\nu}\): specified visible stress tensor.
\item \(\Inv_{33}=T^\D_{\mu\nu}\): residual stress tensor.
\item \(\Inv_{34}=\epsilon_\D=T^\D_{\mu\nu}U^\mu U^\nu/c^2\): residual energy density.
\item \(\Inv_{35}=p_\D=h^{\mu\nu}T^\D_{\mu\nu}/3\): residual isotropic pressure.
\item \(\Inv_{36}=w_\D=p_\D/\epsilon_\D\): residual background equation-of-state ratio.
\item \(\Inv_{37}=h^\alpha{}_\mu T^\D_{\alpha\beta}U^\beta\): residual momentum flux.
\item \(\Inv_{38}=\pi^\D_{\mu\nu}\): residual anisotropic stress.
\item \(\Inv_{39}=T^\D=g^{\mu\nu}T^\D_{\mu\nu}\): residual trace.
\item \(\Inv_{40}=\nabla_\mu T^{\D\,\mu\nu}\): residual divergence diagnostic.
\item \(\Inv_{41}=J^\nu\): interaction current if sectors exchange energy-momentum.
\item \(\Inv_{42}=-\Delta\): scalar Laplacian on \(S^3_A\).
\item \(\Inv_{43}=\lambda_n=n(n+2)/A^2\): scalar eigenvalues.
\item \(\Inv_{44}=g_n=(n+1)^2\): scalar degeneracy.
\item \(\Inv_{45}=\lambda_1=3/A^2\): scalar spectral gap.
\item \(\Inv_{46}=\Xi_{n,A}(x,y)\): scalar addition kernel.
\item \(\Inv_{47}=\Xi_{n,A}(0)=(n+1)^2/(2\pi^2A^3)\): diagonal kernel.
\item \(\Inv_{48}=A_n\ge0\): covariance positivity.
\item \(\Inv_{49}=\sum_nA_n\Xi_{n,A}(0)<\infty\): finite point variance.
\item \(\Inv_{50}=K(s)=\operatorname{Tr}e^{-s\Delta}\): heat trace.
\item \(\Inv_{51}=\zeta_\Delta(s)=\sum_{\lambda>0}\lambda^{-s}\): spectral zeta function.
\item \(\Inv_{52}=\Delta^{X}_{\ell n}\): closed-FLRW transfer function.
\item \(\Inv_{53}=C_\ell^{XY}\): angular power spectrum prediction.
\item \(\Inv_{54}=P_\D(k,z)\) or its discrete analogue: residual clustering spectrum.
\item \(\Inv_{55}=f\sigma_8(z)\): growth observable.
\item \(\Inv_{56}=C_\ell^{\phi\phi}\): CMB lensing spectrum.
\item \(\Inv_{57}=\xi_\pm(\vartheta)\): weak-lensing correlation functions.
\item \(\Inv_{58}={\cal T}_{\rm top}\): topology statistic.
\item \(\Inv_{59}=H_\xi\approx0\): diffeomorphism Hamiltonian constraint on closed slices.
\item \(\Inv_{60}=E_{\BY}(R_A)\): Brown--York quasi-local energy.
\item \(\Inv_{61}=E_{\MS}(R_A)\): Misner--Sharp energy.
\item \(\Inv_{62}=\int_\Sigma\rho\,\dd V=0\): compact Gauss zero-mode condition.
\item \(\Inv_{63}=G_0(x,y)\): zero-subtracted Green function.
\item \(\Inv_{64}=Q_{\partial D}\): quasi-local surface charge on a subregion.
\item \(\Inv_{65}={\cal C}\): data covariance matrix in the likelihood.
\item \(\Inv_{66}=\ln{\cal L}\): predictive likelihood scalar.
\item \(\Inv_{67}=Z_{\cal M}=\int{\cal L}\pi\,\dd\theta\): Bayesian evidence.
\item \(\Inv_{68}=K_{12}=Z_1/Z_2\): Bayes factor between models.
\item \(\Inv_{69}=r_i=d_i-m_i(\theta)\): posterior predictive residual.

\smallskip
\noindent\emph{Conditional mathematical extensions, not required for
the core benchmark.}
\smallskip

\item \(\Inv_{70}=\varphi\): conditional Fibonacci hierarchy eigenvalue.
\item \(\Inv_{71}\): conditional graph adjacency spectrum.
\item \(\Inv_{72}\): conditional quantum phase or Born--Oppenheimer parameter.

\item \(\Inv_{73}=({\cal I},\Theta,\pi,{\cal T},{\cal N},{\cal V},{\cal R})\):
benchmark manifest.
\item \(\Inv_{74}={\cal T}\): invariant-to-observable transfer operator.
\item \(\Inv_{75}={\cal N}\): experimental noise and nuisance model.
\item \(\Inv_{76}={\cal V}\): validation-test vector.
\item \(\Inv_{77}=p(d_{\rm blind}\mid{\cal M})\): blind predictive density.
\item \(\Inv_{78}=r_i^{\rm pp}\): posterior-predictive residual.
\item \(\Inv_{79}=\Delta\chi^2_{\rm pp}\): posterior-predictive discrepancy.
\item \(\Inv_{80}\): versioned reference spectrum table.
\item \(\Inv_{81}\): versioned covariance-kernel table.
\item \(\Inv_{82}\): versioned residual-history table.
\item \(\Inv_{83}\): code-and-data release hash.
\item \(\Inv_{84}\): failure criterion attached to a registered prediction.
\end{enumerate}


\section{Benchmark manifest}
\label{app:manifest}

The following is the minimal manifest for a public benchmark release.
It is deliberately mechanical.
\begin{enumerate}
\item State the geometry: \(A\), \(\Omega_K\), \(a_0\), lapse convention,
      and proper-time derivative.
\item State the visible sector: baryons, photons, neutrinos, and any
      interaction current.  If an interaction current is absent, state
      \(\nabla_\mu T^{\vis\,\mu\nu}=0\).
\item State the residual target: \(T^{\cal D}_{\mu\nu}\),
      \(\epsilon_{\cal D}(a)\), \(p_{\cal D}(a)\), and perturbative
      closure if used.
\item State the spectral sector: scalar, transverse-vector, or TT;
      operator, eigenvalues, degeneracies, and covariance coefficients.
\item State the data vector and covariance exactly as used in the
      likelihood.
\item Publish validation outputs for the analytic limits listed in
      Sec.~\ref{sec:benchmark}.
\item Publish posterior samples, posterior-predictive residuals,
      evidence estimates, and the failure criterion.
\end{enumerate}
This manifest is stricter than a narrative conclusion.  It is also more
valuable: it converts the framework into a reusable standard.

\section{Quarantine rule}
\label{app:quarantine}

The framework excludes unsupported global claims from the PRD/PRA-level
core.  In particular, it does not contain a proof of the Riemann
hypothesis, the Collatz conjecture, Fermat's last theorem, local
connectivity of the Mandelbrot set, or the Poincaré conjecture.  It also
does not use consciousness, reincarnation, or metaphysical claims as
physical assumptions.  Such material may be studied elsewhere as
mathematics, philosophy, or speculation, but it is not part of the
transfer-grade invariant programme unless it supplies definitions,
theorems, observables, and falsifiable likelihoods.

\begin{thebibliography}{99}

\bibitem{Wald1984}
R. M. Wald, \textit{General Relativity} (University of Chicago Press, Chicago, 1984).

\bibitem{HawkingEllis1973}
S. W. Hawking and G. F. R. Ellis, \textit{The Large Scale Structure of Space-Time} (Cambridge University Press, Cambridge, 1973).

\bibitem{Peebles1993}
P. J. E. Peebles, \textit{Principles of Physical Cosmology} (Princeton University Press, Princeton, 1993).

\bibitem{DodelsonSchmidt2020}
S. Dodelson and F. Schmidt, \textit{Modern Cosmology}, 2nd ed. (Academic Press, London, 2020).

\bibitem{Weinberg2008}
S. Weinberg, \textit{Cosmology} (Oxford University Press, Oxford, 2008).

\bibitem{RubinOrdonez1984}
M. A. Rubin and C. R. Ord\'{o}\~nez, Symmetric tensor eigenspectrum of the Laplacian on \(n\)-spheres, J. Math. Phys. \textbf{25}, 2888 (1984).

\bibitem{AbbottSchaefer1986}
L. F. Abbott and R. K. Schaefer, A general, gauge-invariant analysis of the cosmic microwave anisotropy, Astrophys. J. \textbf{308}, 546 (1986).

\bibitem{Jantzen1978}
R. T. Jantzen, Tensor harmonics on the 3-sphere, J. Math. Phys. \textbf{19}, 1163 (1978).

\bibitem{Higuchi1987}
A. Higuchi, Symmetric tensor spherical harmonics on the \(N\)-sphere and their application to the de Sitter group \(SO(N,1)\), J. Math. Phys. \textbf{28}, 1553 (1987); Erratum, J. Math. Phys. \textbf{43}, 6385 (2002).

\bibitem{LindblomTaylorZhang2017}
L. Lindblom, N. W. Taylor, and F. Zhang, Scalar, vector and tensor harmonics on the three-sphere, Gen. Relativ. Gravit. \textbf{49}, 139 (2017), arXiv:1709.08020.

\bibitem{IyerWald1994}
V. Iyer and R. M. Wald, Some properties of Noether charge and a proposal for dynamical black hole entropy, Phys. Rev. D \textbf{50}, 846 (1994), arXiv:gr-qc/9403028.

\bibitem{BrownYork1993}
J. D. Brown and J. W. York, Jr., Quasilocal energy and conserved charges derived from the gravitational action, Phys. Rev. D \textbf{47}, 1407 (1993), arXiv:gr-qc/9209012.

\bibitem{MisnerSharp1964}
C. W. Misner and D. H. Sharp, Relativistic equations for adiabatic, spherically symmetric gravitational collapse, Phys. Rev. \textbf{136}, B571 (1964).

\bibitem{York1972}
J. W. York, Jr., Role of conformal three-geometry in the dynamics of gravitation, Phys. Rev. Lett. \textbf{28}, 1082 (1972).

\bibitem{PageWootters1983}
D. N. Page and W. K. Wootters, Evolution without evolution: Dynamics described by stationary observables, Phys. Rev. D \textbf{27}, 2885 (1983).

\bibitem{Planck2018VI}
N. Aghanim \textit{et al.} (Planck Collaboration), Planck 2018 results. VI. Cosmological parameters, Astron. Astrophys. \textbf{641}, A6 (2020), arXiv:1807.06209.

\bibitem{Planck2018VII}
N. Aghanim \textit{et al.} (Planck Collaboration), Planck 2018 results. VIII. Gravitational lensing, Astron. Astrophys. \textbf{641}, A8 (2020), arXiv:1807.06210.

\bibitem{PlanckPR4CamSpec}
E. Rosenberg, S. Gratton, and G. Efstathiou, CMB power spectra and cosmological parameters from Planck PR4 with CamSpec, Mon. Not. R. Astron. Soc. \textbf{517}, 4620 (2022), arXiv:2205.10869.

\bibitem{DESIDR2BAO2025}
M. Abdul Karim \textit{et al.} (DESI Collaboration), DESI DR2 results II: Measurements of baryon acoustic oscillations and cosmological constraints, arXiv:2503.14738.

\bibitem{DESIDR2Lya2025}
M. Abdul Karim \textit{et al.} (DESI Collaboration), DESI DR2 results I: Baryon acoustic oscillations from the Lyman-alpha forest, arXiv:2503.14739.

\bibitem{DESIDR2DE2025}
K. Lodha \textit{et al.} (DESI Collaboration), Extended dark energy analysis using DESI DR2 BAO measurements, Phys. Rev. D \textbf{112}, 083511 (2025), arXiv:2503.14743.

\bibitem{PantheonPlus2022}
D. Brout \textit{et al.}, The Pantheon+ analysis: Cosmological constraints, Astrophys. J. \textbf{938}, 110 (2022), arXiv:2202.04077.

\bibitem{Moresco2022}
M. Moresco \textit{et al.}, Unveiling the Universe with emerging cosmological probes, Living Rev. Relativ. \textbf{25}, 6 (2022), arXiv:2201.07241.

\bibitem{ACTDR62024}
M. S. Madhavacheril \textit{et al.} (ACT Collaboration), The Atacama Cosmology Telescope: DR6 gravitational lensing map and cosmological parameters, Astrophys. J. \textbf{962}, 113 (2024), arXiv:2304.05203.

\bibitem{CLASS2011}
J. Lesgourgues, The Cosmic Linear Anisotropy Solving System (CLASS) I: Overview, arXiv:1104.2932.

\bibitem{CLASSNonFlat2014}
J. Lesgourgues and T. Tram, Fast and accurate CMB computations in non-flat FLRW universes, JCAP \textbf{09}, 032 (2014), arXiv:1312.2697.

\bibitem{CAMB2000}
A. Lewis, A. Challinor, and A. Lasenby, Efficient computation of cosmic microwave background anisotropies in closed Friedmann-Robertson-Walker models, Astrophys. J. \textbf{538}, 473 (2000), arXiv:astro-ph/9911177.

\bibitem{SahaCompact2025}
S. Saha \textit{et al.} (COMPACT Collaboration), Cosmic topology. Part Ic. Limits on lens spaces from circle searches, JCAP \textbf{01}, 004 (2025), arXiv:2409.02226.

\end{thebibliography}

\end{document}
